\documentclass{report}
\usepackage{amsmath,amssymb}
\setlength{\parindent}{0mm}
\setlength{\parskip}{1em}
\begin{document}
\begin{center}
\rule{6in}{1pt} \
{\large Frederic Nataf \\
{\bf Abstract Robust Coarse Spaces for Systems of PDEs via {Gen}eralized {E}igenproblems in the {O}verlaps}}

Universite P M Curie \\ Laboratoire J L Lions \\ 4 place Jussieu \\ 75005 Paris
\\
{\tt nataf@ann.jussieu.fr}\\
Nicole Spillance\\
Victorita Dolean\\
Patrice Hauret\\
Clemens Pechstein\\
Robert Scheichl\end{center}

{\bf (session on hybrid direct-iterative linear solvers)}

Domain decomposition methods are hybrid direct-iterative parallel
solvers. The effort to reach scalability in domain decomposition methods
has led to the design of so called two-level methods. Each of these
methods is characterized by two ingredients: a coarse space and a
formulation of how this coarse space is incorporated into the domain
decomposition method. We will work in the already extensively studied
framework of the overlapping additive Schwarz preconditioner and focus on
the definition of a suitable coarse space with the aim to achieve
robustness with regard to heterogeneities in any of the coefficients in
the PDEs. This type of problem arises in many applications, such as
subsurface flows or linear elasticity. One way to avoid long stagnation
in Schwarz methods is to build the subdomains in such a way that the
variations in the coefficients are small or nonexistent inside each
subdomain. In this configuration, classical coarse spaces that are based
on these subdomain partitions are known to be robust, see. Recently, some
authors also extended these results for scalar elliptic problems to
certain classes of coefficients that are not resolved by the subdomain
partition and to operator dependent coarse spaces.

However, ideally we would prefer methods that are robust for any
partition into subdomains and regardless of the coefficient distribution.
A class of such stable coarse spaces has been presented among the
literature on Schwarz methods for scalar elliptic PDEs as well as in
earlier work on algebraic multigrid (AMG) method. The key ingredient in
these coarse spaces is the solution of a generalized eigenvalue problem
on each subdomain (or on each coarse element in spectral AMGe). While
these spaces are indeed robust for any arbitrary coefficient
distribution, they do not discriminate between coefficient variations
that influence only the solution in the interior of the subdomains and
those that are actually responsible for the lack of robustness of
standard coarse spaces. A consequence of this is that the resulting
coarse space is often unnecessarily large.\\

Here, we introduce in a variational setting a new coarse space that is
robust and of minimal size. We achieve this by solving local generalized
eigenvalue problems in the overlaps of subdomains that isolate the terms
responsible for slow convergence. The right hand side of the generalized
eigenproblems is constructed from element stiffness matrices and a family
of partition of unity functions/operators. Thus, we only have to assume
access to some topological information to build a suitable partition of
unity and to the element stiffness matrices (as in AMGe methods), which
is reasonable in standard FE packages such as FreeFEM++. The
preconditioners can then be implemented fully algebraically, providing a
viable option for simulations on physical problems without the need to
rewrite entire codes.

We prove a general theoretical result that rigorously establishes the
robustness of the new coarse space and give some numerical examples on
two and three dimensional heterogeneous PDEs and systems of PDEs that
confirm this property.


\end{document}
